\documentclass[11pt]{article}
\usepackage[margin=1in]{geometry}
\usepackage{amsmath,amssymb,amsthm}
\usepackage{physics}
\usepackage{enumitem}
\usepackage{hyperref}
\usepackage{titlesec}

\theoremstyle{definition}
\newtheorem{definition}{Definition}[section]
\newtheorem{example}{Example}[section]
\newtheorem{remark}{Remark}[section]
\newtheorem{proposition}{Proposition}[section]
\newtheorem{theorem}{Theorem}[section]

\titleformat{\section}{\large\bfseries}{Lecture 11.\thesection}{1em}{}

\title{\textbf{Lecture 11: General Measurements} \\ \large Understanding Quantum Information \& Computation}
\author{Based on the lectures by John Watrous (IBM Quantum)}
\date{}

\begin{document}
\maketitle
\tableofcontents
\newpage

%----------------------------------------------------------------------
\section{Overview}
%----------------------------------------------------------------------

This lecture completes the trio of general-description tools by introducing the most general notion of quantum measurement.  We progress through three levels of generality:

\begin{enumerate}
    \item \textbf{Standard basis measurements} (Lecture~1).
    \item \textbf{Projective measurements} (Lecture~3).
    \item \textbf{General measurements / POVMs} (this lecture).
\end{enumerate}

We also discuss \textbf{instruments} (measurements that describe both the outcome probabilities \emph{and} the post-measurement state).

%----------------------------------------------------------------------
\section{Recap: Projective Measurements}
%----------------------------------------------------------------------

A projective measurement is defined by projections $\{\Pi_k\}$ with $\sum_k\Pi_k=I$.  On state $\rho$:
\[
    \Pr(k) = \tr(\Pi_k\,\rho),
    \qquad
    \rho_k = \frac{\Pi_k\,\rho\,\Pi_k}{\tr(\Pi_k\,\rho)}.
\]
Projective measurements are \emph{repeatable}: measuring a second time yields the same outcome.

%----------------------------------------------------------------------
\section{POVMs: Positive Operator-Valued Measures}
%----------------------------------------------------------------------

\begin{definition}[POVM]
A \textbf{POVM} (positive operator-valued measure) is a collection of positive semidefinite operators $\{P_1,\dots,P_m\}$ satisfying
\[
    \sum_{k=1}^{m}P_k = I.
\]
The operators $P_k$ are called \textbf{POVM elements} or \textbf{measurement operators}.
\end{definition}

When a POVM is applied to a state $\rho$, outcome $k$ occurs with probability
\[
    \Pr(k) = \tr(P_k\,\rho).
\]

\begin{remark}
A POVM specifies \emph{only} the outcome probabilities.  It does not, by itself, specify the post-measurement state.  For that, we need an \textbf{instrument} (see below).
\end{remark}

\subsection{Relationship to Projective Measurements}

Every projective measurement is a POVM (with $P_k=\Pi_k$), but not conversely:
\begin{itemize}
    \item POVM elements need not be projections ($P_k^2\neq P_k$ in general).
    \item The number of outcomes $m$ can exceed the dimension of the Hilbert space.
    \item POVM elements need not be orthogonal.
\end{itemize}

\begin{example}[Trine measurement on a qubit]
Three POVM elements, each rank-1, pointing at $120^\circ$ intervals on the Bloch sphere:
\[
    P_k = \frac{2}{3}\ket{\phi_k}\!\bra{\phi_k},\quad k=1,2,3,
\]
where $\ket{\phi_1}=\ket{0}$, $\ket{\phi_2}=\frac{1}{2}(-\ket{0}+\sqrt{3}\ket{1})$, $\ket{\phi_3}=\frac{1}{2}(-\ket{0}-\sqrt{3}\ket{1})$.
Verify: $P_1+P_2+P_3=I$.  This has three outcomes on a two-dimensional system---impossible with a projective measurement.
\end{example}

\subsection{Naimark's Theorem}

\begin{theorem}[Naimark dilation]
Every POVM on a system $X$ can be realised by:
\begin{enumerate}
    \item Appending an ancilla system $A$ initialised to a fixed state.
    \item Applying a joint unitary $V$ on $XA$.
    \item Performing a \emph{projective} measurement on a subsystem.
\end{enumerate}
\end{theorem}

Hence POVMs add no new computational power beyond projective measurements + ancillae, but they are a very convenient abstraction.

%----------------------------------------------------------------------
\section{Quantum Instruments}
%----------------------------------------------------------------------

\begin{definition}[Quantum instrument]
A \textbf{quantum instrument} is a collection of completely positive (CP) maps $\{\Phi_1,\dots,\Phi_m\}$ such that $\sum_k\Phi_k$ is trace-preserving.

When applied to state $\rho$:
\begin{itemize}
    \item Outcome $k$ occurs with probability $\Pr(k)=\tr\bigl(\Phi_k(\rho)\bigr)$.
    \item Post-measurement state given outcome $k$: $\rho_k = \Phi_k(\rho)/\tr\bigl(\Phi_k(\rho)\bigr)$.
\end{itemize}
\end{definition}

The POVM associated with an instrument is $P_k = \sum_j K_{k,j}^\dagger K_{k,j}$ where $\Phi_k(\rho)=\sum_j K_{k,j}\rho\,K_{k,j}^\dagger$.

\begin{example}[Projective measurement as instrument]
$\Phi_k(\rho)=\Pi_k\,\rho\,\Pi_k$.  Single Kraus operator per outcome: $K_k=\Pi_k$.
\end{example}

\begin{example}[L\"uders instrument]
For a POVM $\{P_k\}$, the L\"uders instrument is $\Phi_k(\rho)=\sqrt{P_k}\,\rho\,\sqrt{P_k}$.  This is one of many instruments compatible with a given POVM.
\end{example}

%----------------------------------------------------------------------
\section{Distinguishing Quantum States}
%----------------------------------------------------------------------

\subsection{Minimum-Error State Discrimination}

Given: $\rho_0$ with prior probability $p_0$ and $\rho_1$ with $p_1=1-p_0$.  Design a two-outcome POVM $\{P_0,P_1\}$ to maximise
\[
    p_{\text{correct}} = p_0\tr(P_0\rho_0)+p_1\tr(P_1\rho_1).
\]

\begin{theorem}[Helstrom]
The optimal success probability is
\[
    p_{\text{correct}}^* = \frac{1}{2}\bigl(1+\|p_0\rho_0-p_1\rho_1\|_1\bigr),
\]
where $\|\cdot\|_1$ denotes the trace norm.  The optimal measurement projects onto the positive and negative eigenspaces of $p_0\rho_0-p_1\rho_1$.
\end{theorem}

For pure states $\ket{\psi}$, $\ket{\phi}$ with equal priors:
\[
    p_{\text{correct}}^* = \frac{1}{2}\bigl(1+\sqrt{1-|\!\braket{\psi}{\phi}\!|^2}\bigr).
\]
Perfect discrimination ($p^*=1$) iff $\braket{\psi}{\phi}=0$ (orthogonal), consistent with Lecture~3.

\subsection{Unambiguous State Discrimination}

An alternative paradigm: allow an ``inconclusive'' outcome, but demand zero error when a definite answer is given.  This is possible for non-orthogonal states at the cost of a non-zero inconclusive probability.

%----------------------------------------------------------------------
\section{Relationship Between the Three Measurement Types}
%----------------------------------------------------------------------

\[
    \text{Standard basis} \;\subset\; \text{Projective} \;\subset\; \text{POVM (general)}.
\]

\begin{center}
\begin{tabular}{lccc}
\hline
\textbf{Property} & \textbf{Std.\ basis} & \textbf{Projective} & \textbf{POVM} \\\hline
Elements are projections & Yes & Yes & Not necessarily \\
Elements are orthogonal & Yes & Yes & Not necessarily \\
$\#$outcomes $\le$ dimension & Yes & Yes & Can exceed \\
Repeatable & Yes & Yes & Not necessarily \\
Post-measurement state specified & Yes & Yes & Need instrument \\
\hline
\end{tabular}
\end{center}

%----------------------------------------------------------------------
\section{Summary}
%----------------------------------------------------------------------

\begin{itemize}
    \item A \textbf{POVM} $\{P_k\}$ with $\sum_k P_k=I$ defines the most general quantum measurement (outcome probabilities $\tr(P_k\rho)$).
    \item A \textbf{quantum instrument} additionally specifies the post-measurement state.
    \item \textbf{Naimark's theorem}: every POVM can be implemented with ancillae + a projective measurement.
    \item \textbf{Helstrom's theorem} gives the optimal success probability for distinguishing two states.
    \item POVMs, channels, and density matrices together form the complete general description of quantum information.
\end{itemize}

\end{document}
